%! AP Statistics — Unit 1, Lesson 1.9 — Homework ANSWER KEY
\documentclass[10pt]{article}
\usepackage{apstats-article}
\usepackage{apstats-key}

%=============================================================================
\begin{document}

\pageheader{Unit 1, Lesson 1.9}{Homework --- Answer Key}
\begin{tabularx}{\linewidth}{X X X}
Name:\;\ans{Key} & Date:\; & Period:\;
\end{tabularx}

\vspace{0.10in}

\begin{tcolorbox}[colback=sky, colframe=navy, boxrule=0.8pt, arc=2mm,
    left=3mm, right=3mm, top=1.5mm, bottom=1.5mm]
\small Use explicit comparative language in every sentence that compares groups.
Interpret statistics in context. Show all fence calculations.
\end{tcolorbox}

\vspace{0.12in}

% ════════════════════════════════════════════════════════════════════════════
% PART A
% ════════════════════════════════════════════════════════════════════════════
\begin{blurbbox}[Part A \quad Read and Compare Parallel Boxplots]{navylight}
\small

\begin{center}
\begin{tikzpicture}[scale=1.0]
  \draw[thick] (0.2, 0) -- (7.40, 0);
  \foreach \v/\x in {50/0.00, 55/0.80, 60/1.60, 65/2.40, 70/3.20,
                      75/4.00, 80/4.80, 85/5.60, 90/6.40, 95/7.20} {
    \draw (\x, -0.12) -- (\x, 0.12);
    \node[below, font=\small] at (\x, -0.12) {\v};
  }
  \node[below, font=\small\itshape] at (3.60, -0.56) {Books read per semester};
  \node[left, font=\small] at (0.0, 1.10) {Morning};
  \draw[thick, navy] (1.60, 1.10) -- (2.88, 1.10);
  \draw[thick, navy] (2.88, 0.85) rectangle (4.64, 1.35);
  \draw[thick, navy] (3.68, 0.85) -- (3.68, 1.35);
  \draw[thick, navy] (4.64, 1.10) -- (6.40, 1.10);
  \draw[thick, navy] (1.60, 0.95) -- (1.60, 1.25);
  \draw[thick, navy] (6.40, 0.95) -- (6.40, 1.25);
  \node[left, font=\small] at (0.0, 0.35) {Afternoon};
  \draw[thick, navy] (0.80, 0.35) -- (1.92, 0.35);
  \draw[thick, navy] (1.92, 0.10) rectangle (4.48, 0.60);
  \draw[thick, navy] (3.20, 0.10) -- (3.20, 0.60);
  \draw[thick, navy] (4.48, 0.35) -- (6.08, 0.35);
  \draw[thick, navy] (0.80, 0.20) -- (0.80, 0.50);
  \draw[thick, navy] (6.08, 0.20) -- (6.08, 0.50);
\end{tikzpicture}
\end{center}

\begin{enumerate}[leftmargin=1.5em, itemsep=6pt]

  \item \renewcommand{\arraystretch}{1.5}
        \begin{tabularx}{\linewidth}{@{} l Y Y Y Y Y @{}}
        \rowcolor{sky}
        & $\min$ & $Q_1$ & $M$ & $Q_3$ & $\max$ \\
        \midrule
        Morning   & \ans{60} & \ans{68} & \ans{73} & \ans{79} & \ans{90} \\
        Afternoon & \ans{55} & \ans{62} & \ans{70} & \ans{78} & \ans{88} \\
        \end{tabularx}\\[4pt]
        Morning IQR $= 79 - 68 = $ \ans{11 books}\hfill
        Afternoon IQR $= 78 - 62 = $ \ans{16 books}

  \item \ans{The median number of books read by Morning class students (73) was
        \textit{3 books higher than} the median for Afternoon class students (70),
        suggesting that Morning class students tended to read slightly more per semester.}

  \item \ans{The IQR of books read for the Morning class (11) was \textit{smaller than}
        the IQR for the Afternoon class (16), indicating that Morning class
        reading totals were more consistent from student to student.}

  \item \ans{The two distributions overlap considerably from 62 to 88 books --- a
        range shared by both classes. Large overlap suggests the two classes'
        reading habits are more similar than different, despite the difference
        in medians.}

\end{enumerate}

\end{blurbbox}

\vspace{0.12in}

% ════════════════════════════════════════════════════════════════════════════
% PART B
% ════════════════════════════════════════════════════════════════════════════
\begin{blurbbox}[Part B \quad Back-to-Back Stemplot]{greenacc}
\small

Group X ($n=8$): 62, 65, 67, 71, 74, 76, 81, 83\\
Group Y ($n=8$): 64, 68, 72, 75, 78, 82, 85, 89

\begin{enumerate}[leftmargin=1.5em, itemsep=6pt]

  \item
        \begin{center}
        \small
        \renewcommand{\arraystretch}{1.4}
        \begin{tabular}{r c l}
        \multicolumn{1}{c}{\textit{Group X (reversed)}} & Stem & \textit{Group Y} \\
        \hline
        \ans{7\;\;5\;\;2} & 6 & \ans{4\;\;8} \\
        \ans{6\;\;4\;\;1} & 7 & \ans{2\;\;5\;\;8} \\
        \ans{3\;\;1}      & 8 & \ans{2\;\;5\;\;9} \\
        \hline
        \multicolumn{3}{c}{\small Key: $7 \mid 4 = 74$ for Group X;\enspace $7 \mid 5 = 75$ for Group Y}
        \end{tabular}
        \end{center}

  \item Five-number summary (with $Q_1 = $ avg of 2nd and 3rd values, $M = $ avg of 4th and 5th,
        $Q_3 = $ avg of 6th and 7th for $n=8$):

        \vspace{4pt}
        \renewcommand{\arraystretch}{1.5}
        \begin{tabularx}{\linewidth}{@{} l Y Y Y Y Y @{}}
        \rowcolor{sky}
        & $\min$ & $Q_1$ & $M$ & $Q_3$ & $\max$ \\
        \midrule
        Group X & \ans{62} & \ans{66} & \ans{72.5} & \ans{78.5} & \ans{83} \\
        Group Y & \ans{64} & \ans{70} & \ans{76.5} & \ans{83.5} & \ans{89} \\
        \end{tabularx}

        \vspace{4pt}
        {\small\textit{Group X: $Q_1=(65+67)/2=66$; $M=(71+74)/2=72.5$; $Q_3=(76+81)/2=78.5$.\\
        Group Y: $Q_1=(68+72)/2=70$; $M=(75+78)/2=76.5$; $Q_3=(82+85)/2=83.5$.}}

  \item \ans{\textbf{Shape:} Both distributions are approximately symmetric; the
        stemplot shows leaves spread fairly evenly around the center in each group.\\[4pt]
        \textbf{Center:} The median quiz score for Group Y (76.5) was
        \textit{4 points higher than} the median for Group X (72.5), suggesting
        that Group Y students tended to perform better.\\[4pt]
        \textbf{Spread:} The IQR for Group Y ($83.5 - 70 = 13.5$ points) was
        \textit{slightly larger than} the IQR for Group X ($78.5 - 66 = 12.5$ points),
        meaning Group Y's scores were marginally more spread out in the middle 50\%.
        The range for Group Y ($89 - 64 = 25$ points) also slightly exceeded Group X's
        range ($83 - 62 = 21$ points).}

\end{enumerate}

\end{blurbbox}

\vspace{0.12in}

% ════════════════════════════════════════════════════════════════════════════
% PART C
% ════════════════════════════════════════════════════════════════════════════
\begin{blurbbox}[Part C \quad AP-Style Free Response --- Model Response]{redacc}
\small

\ans{The distribution of books read per semester for the Morning class is
approximately symmetric, with no outliers. The Afternoon class distribution is
also roughly symmetric with no outliers, though slightly more spread out. The
median number of books read by Morning class students (73) was \textit{3 books
higher than} the median for Afternoon class students (70), indicating that
Morning class students tended to read slightly more on average. The IQR of
the Morning class (11 books) was \textit{smaller than} the IQR of the
Afternoon class (16 books), meaning Morning class reading totals were more
consistent from student to student. Despite the difference in spread, the two
distributions overlap considerably (roughly 62 to 88 books), suggesting that
the differences between the classes are modest.}

\end{blurbbox}

\begin{teachernote}
\textbf{Grading notes:}
\begin{itemize}[itemsep=2pt]
  \item Part A: 6 pts. Q1: all 10 values + IQRs correct (2 pts). Q2: center comparison (2 pts: comparative word + both values + context). Q3: spread comparison (1 pt). Q4: overlap + interpretation (1 pt).
  \item Part B: 5 pts. Q5 stemplot (2 pts: all leaves correct). Q6 five-number summaries (2 pts: both rows). Q7 comparison (1 pt: center comparison required; spread credited as bonus).
  \item Part C: 3 pts (1 each for: shape/outlier addressed, center compared with both values, spread compared with both IQRs). All must use comparative language naming both classes.
\end{itemize}
\end{teachernote}

\end{document}
